\section*{第 \ref{ch:b2:plant-plasticity} 章——植物的适应与表型可塑性}

\begin{solution}{exo:b2:plant-plasticity:1}
可塑性：同一\omterm{def:b2:meiosis-heredity:vocabulary}{基因型}依
环境产生多种\omterm{def:b2:meiosis-heredity:vocabulary}{表型}（同一棵橡树上的\omterm{prop:b2:plant-plasticity:sunshade}{阳生叶}与\omterm{prop:b2:plant-plasticity:sunshade}{阴生叶}）。\omterm{def:b2:plant-plasticity:plasticity}{反应规范}：
某一\omterm{def:b2:meiosis-heredity:vocabulary}{基因型}的\omterm{def:b2:meiosis-heredity:vocabulary}{表型}相对于环境
变量（以叶厚随光照的变化为例）。\omterm{def:b2:plant-plasticity:plasticity}{驯化适应}：一种可逆的
生理调节（黑麦在秋季的抗寒锻炼）。
适应：经世代选择的可遗传适合度（
仙人掌的 CAM 代谢）。
\end{solution}

\begin{solution}{exo:b2:plant-plasticity:2}
完整的胚芽鞘从一侧受光：向光弯曲。尖端
以不透明帽罩住：不弯曲——尖端感受
光信号。以透明帽罩住尖端：弯曲——帽本身
无害。基部以不透明套管遮住、尖端暴露：弯曲——
感受在尖端，反应发生在其下方。
\end{solution}

\begin{solution}{exo:b2:plant-plasticity:3}
\omterm{prop:b2:plant-plasticity:sunshade}{阳生叶}：叶小而厚，栅栏组织两到三层，
角质层厚，\omterm{prop:b2:plant-plasticity:stomata}{气孔}多，最大光合速率与
补偿点都高。\omterm{prop:b2:plant-plasticity:sunshade}{阴生叶}：叶大而薄，栅栏组织一层，
单位质量叶绿素更多，呼吸低，补偿点低。
这一选择在叶仍为芽中原基时完成，依据
其所受光照；成熟的叶无法再改变。
\end{solution}

\begin{solution}{exo:b2:plant-plasticity:4}
\omterm{prop:b2:plant-plasticity:stomata}{保卫细胞}吸收钾离子与水分而膨胀，由于其
内壁较厚，向两侧弯开，打开\omterm{prop:b2:plant-plasticity:stomata}{气孔}。蓝光
与胞内低浓度的 $\mathrm{CO_2}$ 在清晨使其开放；\omterm{def:b2:plant-flowering:dormancy}{脱落酸}
（来自干燥土壤中的根或叶片）在干旱时使其关闭，
极干燥的空气同样如此。
\end{solution}

\begin{solution}{exo:b2:plant-plasticity:5}
$\varphi = R/(R + 0.77\,FR)$，当 $FR = 1$ 时：$1.15$：$0.60$；$0.7$：
$0.48$；$0.2$：$0.21$；$0.05$：$0.06$。阈值 $0.4$ 落在
林窗（$0.48$）与冠层（$0.21$）之间：避荫反应
在冠层之下开启。
\end{solution}

\begin{solution}{exo:b2:plant-plasticity:6}
$E = 0.25\times 0.025 = \qty{6.25}{mmol.m^{-2}.s^{-1}}$; $A =
0.25\times 100/1.6 = \qty{15.6}{\micro mol.m^{-2}.s^{-1}}$; $A/E =
2.5\times 10^{-3}$ 摩尔碳每摩尔水，即每个碳对应 400 个水
分子，$400\times 18/12 = \qty{600}{g}$ 水，每
克碳。在 12 小时内：$6.25\times 10^{-3}\times 43200 =
\qty{270}{mol/m^2}$，每平方米 \qty{4.9}{L}。
\end{solution}

\begin{solution}{exo:b2:plant-plasticity:7}
$E = 0.25\times 0.005 = \qty{1.25}{mmol.m^{-2}.s^{-1}}$, $A$
不变：$A/E = 0.0125$，每个碳对应 80 个水分子，\qty{120}{g/g}——
仅为白天代价的五分之一。
\end{solution}

\begin{solution}{exo:b2:plant-plasticity:8}
$v = 2r^{2}\Delta\rho g/9\eta = 2\times 4\times 10^{-12}\times
400\times 9.81/0.09 = \qty{3.5e-7}{m/s}$: \qty{0.35}{\micro m/s},
约 \qty{30}{s} 穿过细胞。在回转器上，细胞相对于重力的
方向在沉降期间每隔几秒改变一次，
因此\omterm{prop:b2:plant-plasticity:gravitropism}{平衡石}从不在一侧聚集，平均
刺激在植物的整合时间（分钟量级）内为零。
\end{solution}

\begin{solution}{exo:b2:plant-plasticity:9}
两小时内平均伸长量 $0.04$；两侧分别为 $1.2\times 0.04 = 0.048$
和 $0.8\times 0.04 = 0.032$；$\theta = 20\times 0.016/2 = 0.16$
弧度，约 \qty{9}{\degree}。弯曲以每小时 $0.08$ 弧度的速率
累积；$\pi/2$ 约需 \qty{20}{h}。
\end{solution}

\begin{solution}{exo:b2:plant-plasticity:10}
贮藏物质可使茎伸长 \qty{50}{cm}。荨麻丛：\qty{30}{cm}
在 7.5 天内耗去 \qty{30}{mg}——幼苗抵达光照时仍有余量
林中：\qty{10}{m} 高度无法企及；幼苗
耗尽其 \qty{50}{mg} 贮藏物质，在十二天内长出半米苍白的茎，
并因黄化而死。若采取耐荫策略——
少量薄叶、缓慢生长——或许可让它存活数年。
\end{solution}

\begin{solution}{exo:b2:plant-plasticity:11}
冰首先在胞外空间形成，那里的溶液
较稀；冰面上的蒸气压低于
细胞内水的蒸气压，水分因而离开细胞加入冰中，细胞
脱水，膜塌陷。经一周低温，细胞
积累糖类与保护性蛋白，以降低其凝固
点，并稳定膜，同时合成抗冻蛋白，
以限制晶体生长。全年持续保持该程序，会将
碳从生长转向保护，并减缓分裂：长期
处于抗寒状态的植物，就是小植株。
\end{solution}

\begin{solution}{exo:b2:plant-plasticity:12}
远红光信号在邻株可能被超越时可信，
若邻株是树木则是误导；暖和的一周在四月可信，
二月则是误导。在密植作物中，每株植物的邻株都是其
同类，因此信号形式上可信却无用，育种者
便选育忽略它的植物。\omterm{def:b2:plant-plasticity:plasticity}{反应规范}是记录
某一谱系过去各线索说真话的频率，它
在现状不同于过去时即失效。
\end{solution}

\begin{solution}{pb:b2:plant-plasticity:1}
\textbf{1.} 开阔处：$1.15/(1.15 + 0.77) = 0.60$；冠层：$0.2/(0.2 +
0.77) = 0.21$。
\textbf{2.} 低 $\varphi$：头顶有植物遮挡——伸长。在
遮阳网下 $\varphi$ 保持 $0.60$：这是阴天，而非邻株；没有
避荫反应，只是因光照不足而生长变慢。
\textbf{3.} $1 + 2(0.60 - 0.21) = 1.78$：每天 \qty{2.7}{cm}。
\textbf{4.} \qty{37}{cm} 对比 \qty{21}{cm}。
\textbf{5.} 向林窗弯曲（\omterm{prop:b2:plant-plasticity:phototropism}{向光性}，经蓝光
受体向光素），以及较快的伸长和叶片重新定向
在 $R{:}FR$ 较高的一侧（经\omterm{prop:b2:plant-flowering:photoperiod}{光敏色素}）。
\textbf{6.} \omterm{prop:b2:plant-plasticity:sunshade}{阴生叶}做得大而薄，以便在稀少的
光照下低成本捕光；在全日照下它们会过热、发生光抑制并灼伤，
植物必须以\omterm{prop:b2:plant-plasticity:sunshade}{阳生叶}取而代之。
\textbf{7.} 在 \qty{15}{mm} 的区段上每天 \qty{2.7}{cm}，即
每小时 \qty{0.11}{cm}，相对伸长率为每小时 $0.074$；
背荫侧 $1.3\times 0.074 = 0.096$，向光侧 $0.7\times 0.074 =
0.052$。
\textbf{8.} $\theta = 15\times(0.096 - 0.052)/2 = 0.33$ 弧度，约
\qty{19}{\degree}，在一小时内。
\textbf{9.} \qty{60}{\degree} 即 $1.05$ 弧度：约三小时。
\textbf{10.} $v = 2\times 6.25\times 10^{-12}\times 400\times
9.81/0.09 = \qty{5.5e-7}{m/s}$, \qty{0.55}{\micro m/s}: \qty{22}{s}
穿过细胞。
\textbf{11.} \omterm{prop:b2:plant-plasticity:phototropism}{向光性}；避荫反应本身提高了
枝条愿意生长的角度：光是
争夺的资源，而逆重力伸直的茎将
重新转向荫蔽之处。
\textbf{12.} 根向下弯曲，直至竖直：\omterm{prop:b2:plant-plasticity:gravitropism}{平衡石}
沉降到下侧，\omterm{prop:b2:plant-meristems:apical}{生长素}在那里积累，而在根中
额外的\omterm{prop:b2:plant-meristems:apical}{生长素}抑制生长，因此上侧生长更快。同样的
重新分布，但细胞的敏感性相反。
\textbf{13.} $E = 0.2\times 0.015 = \qty{3}{mmol.m^{-2}.s^{-1}}$；在
$2\times 10^{-3}$ \unit{m^2} 与 \qty{50400}{s} 内：\qty{0.30}{mol}，
每天 \qty{5.4}{g} 水。
\textbf{14.} $A = 0.2\times 80/1.6 = \qty{10}{\micro mol.m^{-2}.s^{-1}}$;
每天每平方米 $0.50$ 摩尔碳；$A/E = 3.3\times
10^{-3}$：每个碳对应 300 个水分子，\qty{450}{g/g}。
\textbf{15.} $0.50\times 2\times 10^{-3} = \qty{1}{mmol}$: \qty{12}{mg}
每天的碳量。
\textbf{16.} 林窗：$E = \qty{6}{mmol.m^{-2}.s^{-1}}$，\qty{10.9}{g}
每天；效率 \qty{900}{g/g}。
\textbf{17.} $E = 0.02\times 0.03 = \qty{0.6}{mmol.m^{-2}.s^{-1}}$
（降为十分之一）；$A = 0.02\times 200/1.6 = \qty{2.5}{\micro
mol.m^{-2}.s^{-1}}$（降为四分之一）。水分下降更多：叶片
降低其内部 $\mathrm{CO_2}$ 时，碳的梯度
扩大，而水分的梯度则不能。
\textbf{18.} 含水 \qty{4.25}{g}，可失去 \qty{0.85}{g}；按
每天 \qty{10.9}{g}（每小时 $0.78$ 克），约一小时即萎蔫。
\textbf{19.} CAM 缓慢，需要肉质贮藏组织；
在冠层下争夺光照的幼苗，失水本来就很少，
因此需要的是速度，而非节约。
\textbf{20.} $100/2 = \qty{50}{cm}$.
\textbf{21.} 荨麻丛：15 天内长 \qty{40}{cm}，耗去 \qty{80}{mg}——
它便能抵达光照。林中：\qty{8}{m} 高，够不着；它耗尽
贮藏物质于半米长的茎上并死亡。在林中更优的做法是：保持矮小，
长出薄的\omterm{prop:b2:plant-plasticity:sunshade}{阴生叶}，等待林窗——这是山毛榉的做法。
\textbf{22.} 桦树：陡峭的\omterm{def:b2:plant-plasticity:plasticity}{反应规范}，在低
$\varphi$ 时强烈伸长，因为先锋植物的邻株与其大小相当。山毛榉：
平缓的\omterm{def:b2:plant-plasticity:plasticity}{反应规范}，伸长少而耐受性强，因为其邻株
是树木。
\textbf{23.} 关闭避荫反应——保持
\omterm{prop:b2:plant-flowering:photoperiod}{光敏色素} B 保持活性，或去除它所
抑制的\omterm{prop:b2:cell-differentiation:transcriptional}{转录因子}——并检查开花时间、萌发和茎的
强度，这些均受同一途径调控。
\textbf{24.} 不能感知邻株的植物，在本可取胜时
便会被超越；总是相信线索的植物，
会在每棵树下白白耗费自己。胜出的\omterm{def:b2:plant-plasticity:plasticity}{反应规范}是押注于
线索，押注程度与该线索在谱系过去中
正确的次数成比例。
\textbf{25.} 开阔处 $\varphi = 0.60$，冠层之下 $0.21$；
\qty{37}{cm}（14 天后）；阴处每天失水 \qty{5.4}{g}，
林窗中 \qty{10.9}{g}。
\end{solution}
